% !TeX root = ../../tesi.tex

\section{Obiettivi e vincoli del progetto}
\label{sec:obiettivi}

Mi sono candidato a luglio 2026 e ad agosto ho sostenuto un colloquio conoscitivo con il
\emph{tutor} e con il fondatore. Halue mi ha poi proposto due temi, l’interfaccia generativa e la
gestione del contesto di un agente. Ho scelto il primo e ho proposto di declinarlo in un assistente
che seleziona e popola componenti dell’interfaccia con i dati del negozio.

Il piano organizza il prototipo su due livelli (\figref{fig:livelli}). L’\emph{app layer} è il
\emph{widget} di chat nel negozio e si limita a visualizzare la conversazione e i componenti.
L’\emph{agent layer} è il \emph{backend} che ospita l’agente e conserva lo stato della
conversazione. Per interrogare le API di Shopify l’agente usa il
\gloss{tool-calling}{\emph{tool calling}}, il meccanismo con cui un modello linguistico chiede al
programma che lo ospita di eseguire una funzione e ne riceve il risultato.

\begin{figure}[htbp]
	\centering
	\setstretch{1}
	% L'architettura su due livelli prevista dal piano di lavoro: il widget nel negozio
% (app layer) e il backend che ospita l'agente (agent layer), con i servizi esterni a destra.
% Il disegno è largo quanto il testo (150 mm): si include con \input, senza ridimensionarlo.

\begin{tikzpicture}[
	every node/.style={text=FiguraTesto},
	blocco/.style={draw=FiguraTenue, fill=white, rounded corners=3pt, text width=#1-2mm,
		minimum width=#1, minimum height=13mm, align=center, inner sep=1.5pt,
		font=\fontsize{7.5}{9}\selectfont,
		execute at begin node={\hyphenpenalty=10000\exhyphenpenalty=10000}},
	blocco/.default=32mm,
	etichetta/.style={font=\fontsize{7}{8.5}\selectfont, text=FiguraTenue, inner sep=1pt,
		align=center},
]
	\begin{scope}[on background layer]
		\path[partizione] (0mm, 0mm) rectangle (38mm, -44mm);
		\path[intestazione] (0mm, 0mm) rectangle (38mm, -6mm);
		\node[titolo] at (19mm, -3mm) {\emph{App layer}: negozio};
		\path[partizione] (72mm, 0mm) rectangle (110mm, -44mm);
		\path[intestazione] (72mm, 0mm) rectangle (110mm, -6mm);
		\node[titolo] at (91mm, -3mm) {\emph{Agent layer}: \emph{backend}};
		\path[partizione] (122mm, 0mm) rectangle (150mm, -44mm);
		\path[intestazione] (122mm, 0mm) rectangle (150mm, -6mm);
		\node[titolo] at (136mm, -3mm) {Servizi esterni};
	\end{scope}

	\node[blocco] (widget) at (19mm, -16mm)
		{\voce{\emph{Widget} di chat}{mostra la conversazione}};
	\node[blocco] (componenti) at (19mm, -34mm)
		{\voce{Libreria di componenti}{popolati con i dati del negozio}};
	\node[blocco] (agente) at (91mm, -16mm)
		{\voce{Agente}{sceglie testo o componente, esegue gli strumenti}};
	\node[blocco] (stato) at (91mm, -34mm)
		{\voce{Stato}{storico della conversazione}};
	\node[blocco=24mm] (llm) at (136mm, -16mm) {\voce{Modello linguistico}{genera le risposte}};
	\node[blocco=24mm] (api) at (136mm, -34mm) {\voce{API di Shopify}{catalogo, carrello, ordini}};

	\draw[freccia] ([yshift=3mm]widget.east) -- node[etichetta, above=0.5mm]
		{messaggio del cliente} ([yshift=3mm]agente.west);
	\draw[freccia] ([yshift=-3mm]agente.west) -- node[etichetta, below=0.5mm]
		{testo o componente\\e tre messaggi suggeriti} ([yshift=-3mm]widget.east);
	\draw[freccia] (widget) -- node[etichetta, right=0.5mm] {disegna} (componenti);
	\draw[freccia, {Stealth[length=1.8mm,width=1.4mm]}-{Stealth[length=1.8mm,width=1.4mm]}]
		(agente) -- (stato);
	\draw[freccia, {Stealth[length=1.8mm,width=1.4mm]}-{Stealth[length=1.8mm,width=1.4mm]}]
		(agente.east) -- (llm.west);
	\draw[freccia] (agente.east) ++(0, -3mm) -- ++(6mm, 0) |- (api.west);
\end{tikzpicture}

	\caption[L’architettura su due livelli prevista dal piano di lavoro]{L’architettura su due
	livelli prevista dal piano di lavoro. Il \emph{widget} non contiene logica conversazionale: ogni
	decisione spetta all’agente, che può interrogare più volte modello e API prima di rispondere.}
	\label{fig:livelli}
\end{figure}

Gli obiettivi hanno un identificativo, che uso nel resto della relazione, e tre livelli di
priorità. Gli obiettivi obbligatori (\tabref{tbl:obbligatori}) sono vincolanti e, oltre a
dimostrare la fattibilità dell’approccio, chiedono di misurarne l’efficacia.

Gli obiettivi desiderabili hanno un valore aggiunto riconoscibile, quelli facoltativi un valore
aggiunto non prioritario (\tabref{tbl:desiderabili}). Riprendono tre aspetti che il \emph{tutor}
ha indicato come interessanti per l’azienda: il supporto a più mercati, i temi Hydrogen, in cui il
riuso dei componenti React è più semplice e l’integrazione con il carrello più complessa, e una
ricerca nel catalogo scritta su misura, perché il \emph{tutor} giudicava elementare quella del
\emph{server} MCP di Shopify.

\begin{tabellatesi}{P{0.1}P{0.9}}
	{Gli obiettivi obbligatori del piano di lavoro.}
	{tbl:obbligatori}
	\toprule
	\rowcolor{TabellaIntestazione} \textbf{ID} & \textbf{Obiettivo} \\
	\midrule
	OB1 & \emph{App layer}: \emph{widget} di chat in React e TypeScript integrato nel negozio \\
	OB2 & \emph{Agent layer}: \emph{backend} Node.js con un agente dotato di \emph{tool calling}, in
	grado di interrogare le API di Shopify \\
	OB3 & Risposte in linguaggio naturale (domande generali, sui prodotti, di chiarimento) oppure con
	un componente predefinito, scelto e popolato in base alla conversazione \\
	OB4 & Libreria di almeno due o tre componenti flessibili e tre messaggi suggeriti per ogni
	risposta \\
	OB5 & Raccolta di metriche di coinvolgimento e di conversione sull’uso del \emph{widget} \\
	\bottomrule
\end{tabellatesi}

\begin{tabellatesi}{P{0.1}P{0.2}P{0.7}}
	{Gli obiettivi desiderabili e facoltativi del piano di lavoro.}
	{tbl:desiderabili}
	\toprule
	\rowcolor{TabellaIntestazione} \textbf{ID} & \textbf{Priorità} & \textbf{Obiettivo} \\
	\midrule
	DE1 & desiderabile & Supporto a \gloss{shopify-markets}{Shopify Markets}, la funzione con cui un
	negozio vende in più paesi con lingue e valute diverse \\
	DE2 & desiderabile & Componente per lo stato dell’ordine e il supporto dopo la vendita \\
	FA1 & facoltativo & Prova di fattibilità dell’integrazione su Hydrogen, con il riuso dei
	componenti del \emph{widget} \\
	FA2 & facoltativo & Confronto sperimentale tra la ricerca su misura, il \emph{server} MCP di
	Shopify e un eventuale \emph{server} MCP proprio \\
	\bottomrule
\end{tabellatesi}

Il piano ripartisce 310 ore indicative tra dieci attività (\figref{fig:piano}) e prevede che la
ripartizione sia rifinita con il \emph{tutor} nelle prime settimane. Nel capitolo 3 la confronto
con l’uso effettivo del tempo.

\begin{figure}[htbp]
	\centering
	\setstretch{1}
	% La ripartizione indicativa delle 310 ore prevista dal piano di lavoro, come grafico a torta.
% Le fette seguono l'ordine del piano in senso orario dall'alto: in grigio chiaro le attività che
% precedono il codice, in bordeaux lo sviluppo dei due livelli, in grigio medio quelle successive.
% Il disegno è largo quanto il testo (150 mm): si include con \input, senza ridimensionarlo.

\begin{tikzpicture}[every node/.style={text=FiguraTesto}]
	\def\raggio{22mm}
	\def\inizio{90}
	% ore/colore/lato (1 a destra, -1 a sinistra)/quota dell'etichetta/titolo/dettaglio
	\foreach \ore/\colore/\lato/\quota/\titolo/\dettaglio in {
		20/TabellaIntestazione/1/27/{Studio}/{Shopify e strumenti per agenti},
		20/TabellaIntestazione/1/16/{Analisi dei requisiti}/{e perimetro del progetto},
		20/TabellaIntestazione/1/5/{Progettazione}/{dell’architettura},
		40/TabellaIntestazione/1/-7/{Negozio di sviluppo}/{e accesso ai dati},
		50/FiguraAccento/1/-21/{\emph{Agent layer}}/{sviluppo del \emph{backend}},
		60/FiguraAccento!65/-1/-21/{\emph{App layer}}/{sviluppo del \emph{widget}},
		30/FiguraBordo/-1/-7/{Integrazione}/{del flusso conversazionale},
		30/FiguraBordo/-1/5/{Metriche}/{raccolta dei dati di utilizzo},
		20/FiguraBordo/-1/16/{\emph{Test}}/{flusso utente e casi limite},
		20/FiguraBordo/-1/27/{Documentazione}/{tecnica}} {
		\pgfmathsetmacro\fine{\inizio - \ore/310*360}
		\pgfmathsetmacro\medio{(\inizio + \fine)/2}
		\fill[\colore, draw=white, line width=0.8pt] (0, 0) -- (\inizio:\raggio)
			arc[start angle=\inizio, end angle=\fine, radius=\raggio] -- cycle;
		\draw[FiguraTenue] (\medio:\raggio + 1mm) -- (\lato*\raggio + \lato*10mm, \quota mm)
			-- ++(\lato*3mm, 0);
		\ifnum\lato=1
			\node[anchor=west, align=left, font=\fontsize{7.5}{9}\selectfont]
				at (\raggio + 14mm, \quota mm) {\voce{\titolo, \ore{} ore}{\dettaglio}};
		\else
			\node[anchor=east, align=right, font=\fontsize{7.5}{9}\selectfont]
				at (-\raggio - 14mm, \quota mm) {\voce{\titolo, \ore{} ore}{\dettaglio}};
		\fi
		\global\let\inizio\fine
	}
\end{tikzpicture}

	\caption[La ripartizione delle ore prevista dal piano di lavoro]{La ripartizione delle ore
	prevista dal piano di lavoro, in senso orario dall’alto. In grigio chiaro le attività che
	precedono la scrittura del codice, in bordeaux lo sviluppo dei due livelli, in grigio medio le
	attività che lo completano.}
	\label{fig:piano}
\end{figure}

Per l’azienda il criterio di successo era un prototipo funzionante dall’inizio alla fine su un
negozio di esempio, accompagnato da una valutazione onesta di ciò che funziona e di ciò che non
funziona. Valuto il grado di soddisfacimento di ogni obiettivo nella \secref{sec:soddisfacimento}.

Il progetto aveva i seguenti vincoli:

\begin{itemize}
	\item \textbf{Quantità di lavoro}: da 300 a 320 ore, con una distribuzione al massimo a tempo
	pieno, dal 14 settembre al 9 novembre 2026.
	\item \textbf{Luogo di lavoro}: da remoto, con un giorno a settimana in sede, un allineamento
	quotidiano con il \emph{tutor} e una riunione settimanale di un’ora.
	\item \textbf{Tecnologie}: TypeScript in tutto il progetto. Il piano rimandava alle prime settimane la scelta degli strumenti
	per costruire l’agente e del modello linguistico, che il \emph{tutor} ha lasciato a me.
	\item \textbf{Piattaforma}: prima i temi Liquid, che sono i più diffusi, e Hydrogen soltanto come
	eventuale approfondimento finale.
	\item \textbf{Ambiente}: un unico negozio di sviluppo, senza clienti né traffico reali, con due
	temi installati per le prove. L’ho allestito come un negozio di cosmetica botanica con circa
	quaranta prodotti, perché il comportamento di un assistente dipende dal catalogo su cui opera, e
	varianti, collezioni, confezioni regalo, prodotti esauriti e archiviati lo avvicinano a un caso
	reale.
	\item \textbf{Natura del prodotto}: un prototipo di ricerca e sviluppo, essenziale ma solido,
	senza i requisiti di scala e di costo di un prodotto in produzione.
\end{itemize}
